\section{Introduction}
\label{sec:introduction}

Flooding and desertification are visible manifestations of
environmental stress that can affect transportation, infrastructure,
agriculture, ecosystems, and human settlements. Aerial and
remote-sensing imagery provides an efficient mechanism for monitoring
large geographic regions; however, automatic scene classification
remains difficult when imagery is collected from heterogeneous
platforms, locations, resolutions, and source datasets.

Deep convolutional neural networks can achieve high classification
accuracy when training and test images are sampled from similar
distributions. Such performance may not reflect deployment robustness,
particularly when images from the same source characteristics are
present in both training and test partitions.

This study therefore treats source-aware evaluation and domain shift
as central components of climate-scene classification. The task is
three-class classification of aerial images into \emph{Flooded},
\emph{Desert}, and \emph{Neither}. The primary dataset contains
6,334 images assembled from multiple underlying sources.

The investigation addresses the following research questions:

\begin{itemize}
    \item \textbf{RQ1:} How much does performance change between
    conventional mixed evaluation and source-disjoint evaluation?

    \item \textbf{RQ2:} How does a MobileNetV2 ImageNet
    transfer-learning pipeline compare with a custom CNN under
    cross-source testing?

    \item \textbf{RQ3:} Do conservative domain-generalization
    augmentation and Representation Self-Challenging improve
    robustness consistently across random seeds?

    \item \textbf{RQ4:} Do the observed failure patterns persist on an
    independent external UAV dataset?

    \item \textbf{RQ5:} Which visual and scene-composition cues are
    associated with systematic external false-flood predictions?
\end{itemize}

The main contributions of this work are:

\begin{enumerate}
    \item high-confidence per-image reconstruction of source
    provenance for the 6,334-image climate dataset;

    \item mixed and source-disjoint evaluation protocols designed to
    expose cross-source generalization failure;

    \item multi-seed comparison of ERM, DG-AUG-V2, and RSC;

    \item external UAV validation using FloodNet with leakage
    screening;

    \item frozen Grad-CAM analysis of correct and systematic
    false-flood predictions; and

    \item post-hoc pixel-level segmentation analysis connecting
    consensus false-flood predictions with ordinary water,
    non-flooded roads, and heterogeneous scene composition.
\end{enumerate}

% TODO: add verified literature citations here.
